\section*{結言}

この24時間のXを通して見えるのは、生成AI競争の評価軸が急速に多層化していることである。Claudeのウォーターマークでは規制と生成物識別、GrokやGeminiではエージェントやcomputer-use、DeepSeekやQwenでは価格とローカル実行が話題になった。単純な「最も賢いモデル」を決める競争から、用途ごとに速度、価格、運用性、制御性を比較する段階へ移っている。

また、公式発表そのものが少ない時間帯でも、Xではデモ動画、ベンチマーク画像、価格情報、開発者や経営者の短い発言を通じて評価が継続的に更新される。とりわけ中国系モデルと米国系frontier modelが同じ投稿内で比較される場面が増え、利用者にとってモデルの出自よりも「手元のタスクでどう動くか」が重要になっている。

Daily Xはこうした速報性の高いコミュニティの温度を記録する。個々の性能主張や価格、一次資料との整合性は週刊版で改めて検証するとして、8月16日朝時点では、モデル単体の更新よりも、実利用における比較とエージェント運用の成熟が一日の主要テーマだった。
